% No numbered results (statements-of-record rows resolved here: none).
\section{Conclusion}
\label{sec:conclusion}

The Riemann hypothesis is the vanishing of the positive ledger
\(\AZ=\sum_\rho m_\rho(\Real\rho-\frac12)^2\) of squared displacements of
the nontrivial zeros from the fixed line of the involution
\(s\mapsto1-\bar s\).  If a recognition source supplies domination
records of vanishing trace together with a bridge encoding these displacements in an auxiliary
readout, the duality-confinement theorem places every represented zero on the
critical line of its coordinates, and, with a map from coordinates to the reals that sends \(\frac12\) to
\(\frac12\) and carries the represented zeros onto the real parts of all
nontrivial zeros, the Riemann hypothesis follows.  On any shell this
hypothesis is equivalent to the critical-line statement for the
represented zeros, hence on the bare shell to the Riemann hypothesis, and so
is its concrete form on finite windows closed under the reflection
\(s\mapsto1-\bar s\), with unit weights, that eventually contain every
zero.  For windows closed under \(s\mapsto1-s\), with invariant weights, an
operator-level transport bound for the full identity target cannot
tend to zero on a nonempty window, and several natural forms of a
payment for the window sums along windows that eventually contain
every zero are again equivalent to the Riemann hypothesis.  These facts are elementary; they use analytic properties of \(\zeta\)
such as the functional equation and the zero-free half-plane
\(\Real s\geq1\), but no quantitative or prime-side input.  The paper thereby
describes what a proof along this route would need, shows that the
obvious candidates are no easier than the Riemann hypothesis, and does
not claim to supply one.  The mathematical statements are formalized in Lean~4 with Mathlib
(\Cref{app:formalization}).
